\textbf{1. Kiến thức cần nhớ}
\begin{itemize}
\item Với $A$, $B$ là các biểu thức tùy ý:
\begin{itemize}
\item Tổng hai lập phương: $\boxed{A^3+B^3=(A+B)\left(A^2-AB+B^2\right)}$.
\item Hiệu hai lập phương: $\boxed{A^3-B^3=(A-B)\left(A^2+AB+B^2\right)}$.
\end{itemize}
\item $A^2-AB+B^2$ và $A^2+AB+B^2$ gọi là \textit{bình phương thiếu} của hiệu và của tổng (hạng tử giữa là $AB$, không phải $2AB$).
\item Phân biệt: $(A+B)^3$ là \textit{lập phương của một tổng}; $A^3+B^3$ là \textit{tổng hai lập phương}.
\end{itemize}

\textbf{2. Phương pháp giải các dạng}

%% DANG 8D2?8-1
\textbf{Dạng 1. Khai triển, viết dưới dạng tích}
\begin{enumerate}
\item Viết mỗi hạng tử thành lập phương: $8x^3=(2x)^3$, $27=3^3$, $\dfrac18=\left(\dfrac12\right)^3$.
\item Áp dụng $A^3\pm B^3=(A\pm B)\left(A^2\mp AB+B^2\right)$.
\item Ngược lại, gặp tích $(A\pm B)(A^2\mp AB+B^2)$ thì viết gọn thành $A^3\pm B^3$.
\end{enumerate}
Lưu ý: dấu trong ngoặc thứ hai ngược với dấu trong ngoặc thứ nhất.

%% DANG 8D2?8-2
\textbf{Dạng 2. Rút gọn, tính giá trị biểu thức}
\begin{enumerate}
\item Nhận ra các tích có dạng $(A\pm B)(A^2\mp AB+B^2)$, viết gọn thành $A^3\pm B^3$.
\item Kết hợp các hằng đẳng thức khác (lập phương, bình phương), thu gọn.
\item Thay giá trị của biến vào biểu thức đã rút gọn; có thể dùng $A^3+B^3=(A+B)^3-3AB(A+B)$ khi biết $A+B$, $AB$.
\end{enumerate}
